%=====================================================================
% 光学（02）Lecture 4 · 作业答案卡（English，供手抄上交）
% 来源：Electro-optic & Magneto-optic Effects.pdf p82（= Homework EO MO.docx）
% 数值核对：subjects/02-光学/讲义/eomo-check/verify_l4.py
% 编译：xelatex 两遍
%=====================================================================
\documentclass[10.5pt]{ctexart}
%=====================================================================
% theme-violet.tex —— 讲义精讲 · 紫花主题
% 用法：\documentclass[11pt]{ctexart}
%       \input{assets/theme-violet}
%       \begin{document} ... \end{document}
% 只用 ctexart 自带能力 + xcolor/enumitem/fancyhdr/pifont/amsmath，
% 不依赖 tikz/tcolorbox，保证 MiKTeX 可直接编译。
%=====================================================================

\usepackage[a4paper,left=2.0cm,right=2.0cm,top=2.1cm,bottom=2.3cm,%
            headsep=12pt,footskip=22pt]{geometry}
\usepackage{amsmath,amssymb}
\usepackage{needspace}
\usepackage{xcolor}
\usepackage{array}
\usepackage{enumitem}
\usepackage{pifont}
\usepackage{fancyhdr}
\usepackage[hidelinks]{hyperref}

%---------------------------------------------------------------- 配色
\definecolor{violet}{RGB}{107,78,143}      % 主色：紫
\definecolor{violetdeep}{RGB}{63,44,92}    % 深紫：标题/强调
\definecolor{petal}{RGB}{183,159,208}      % 花瓣：装饰线、花饰
\definecolor{bandviolet}{RGB}{247,243,251} % 浅紫底：档案/提示框
\definecolor{bandcream}{RGB}{252,249,242}  % 米白底：易错/注意框
\definecolor{rulepurple}{RGB}{217,207,230} % 分隔线
\definecolor{inkgray}{RGB}{92,96,110}      % 次要文字

%---------------------------------------------------------------- 字体
% 有 Noto CJK 用 Noto，没有就退回 ctex 默认（Windows 下为宋体/黑体）
\IfFontExistsTF{Noto Serif CJK SC}{\setCJKmainfont{Noto Serif CJK SC}}{}
\IfFontExistsTF{Noto Sans CJK SC}{\setCJKsansfont{Noto Sans CJK SC}}{}

%---------------------------------------------------------------- 版式
\setlist{nosep,leftmargin=1.9em}
\linespread{1.28}
\setlength{\parskip}{3.5pt}
\setlength{\parindent}{0pt}
\setlength{\tabcolsep}{5pt}
\renewcommand{\arraystretch}{1.28}

%---------------------------------------------------------------- 花饰
% \ding{95} = 五瓣小花，\ding{96} = 星形花；都来自 pifont，必定存在
\newcommand{\orn}{\textcolor{petal}{\ding{95}}}          % 行内小标
\newcommand{\ornstar}{\textcolor{petal}{\ding{96}}}      % 星形花
\newcommand{\flowerrule}{%                                % 居中花饰分隔线
  \par\vspace{7pt}%
  \noindent\makebox[\linewidth]{\color{petal}%
    \rule{0.20\linewidth}{0.6pt}\hspace{8pt}{\small\ding{95}}\hspace{4pt}%
    \raisebox{1.1pt}{\tiny\ding{95}}\hspace{4pt}{\small\ding{95}}\hspace{8pt}%
    \rule{0.20\linewidth}{0.6pt}}%
  \par\vspace{7pt}}

%---------------------------------------------------------------- 标题与节头
\newcommand{\doctitle}[3]{%   #1 主标题  #2 副标题  #3 元信息（可用 \\ 分行）
  \begin{center}
    {\color{petal}\small \orn\enspace #2\enspace \ornstar}\\[7pt]
    {\LARGE\bfseries\color{violetdeep} #1}\\[10pt]
    {\color{petal}\rule{0.55\linewidth}{1pt}}\\[4pt]
    {\small\color{inkgray} #3}
  \end{center}\vspace{2pt}}

\newcommand{\hdr}[2]{%   #1 章节号（如 §0）  #2 标题
  \par\vspace{12pt}%
  \noindent\colorbox{violetdeep}{\color{white}\bfseries\small\;#1\;}\hspace{7pt}%
  {\color{violetdeep}\large\bfseries #2}\par\nobreak
  \vspace{3pt}{\color{petal}\hrule height 1.2pt}\vspace{6pt}}

\newcommand{\subhdr}[1]{%
  \par\vspace{6pt}%
  \noindent{\color{violet}\bfseries #1}\par\vspace{1pt}}

%---------------------------------------------------------------- 提示框
\newcommand{\callout}[1]{% 浅紫底：讲义档案、读法说明
  \par\medskip\noindent
  \fcolorbox{rulepurple}{bandviolet}{\parbox{\dimexpr\linewidth-2\fboxsep-2\fboxrule}{\vspace{2pt}#1\vspace{2pt}}}%
  \par\medskip}

\newcommand{\warnbox}[1]{% 米白底：易错、边界条件
  \par\medskip\noindent
  \fcolorbox{rulepurple}{bandcream}{\parbox{\dimexpr\linewidth-2\fboxsep-2\fboxrule}{\vspace{2pt}#1\vspace{2pt}}}%
  \par\medskip}

\newcommand{\gapitem}[2]{% #1 缺口类型（紫） #2 内容
  \par\vspace{2pt}\noindent
  \hangindent=1.4em\hangafter=1
  {\color{petal}\small\ding{96}}\enspace{\bfseries\color{violetdeep}#1}\enspace #2\par}

% 题干块（\qstem）：每道例题／作业的「详解」开头必须先复述题干与出处。
% 起因：2026-10-25 第五章讲义把七道作业的题面集中放在一页、详解另起一页，
% 读者翻到「第 4 题」时看不到条件，等于一道没有题干的题目。
% 规则：详解块内自带题干，读者不翻页就能读懂题在问什么。
\newcommand{\qstem}[1]{%
  \par\medskip\noindent
  \fcolorbox{violet}{bandviolet}{%
    \parbox{\dimexpr\linewidth-2\fboxsep-2\fboxrule}{%
      \vspace{2pt}%
      {\small\bfseries\color{violet}题干}\hspace{0.8em}#1%
      \vspace{2pt}}}%
  \par\medskip}

% v2: 保护题干与开头，长题可跨页，续页补上下文。
\newcommand{\qbreak}{\needspace{6\baselineskip}}
% 独立练习：ID、来源、题干、考点、答案标签（通常 ans:ID）。
\newcommand{\practice}[5]{%
  \qbreak\subhdr{练习 #1}\qstem{（#2）#3}%
  \textbf{考点}：#4\par
  详解见附录第~\pageref{#5}~页。\par\medskip}
% 附录开头：ID、来源、条件摘要；自动产生 ans:ID 标签。
\newcommand{\answerstart}[3]{%
  \qbreak\subhdr{#1\quad 详解}\label{ans:#1}\qstem{（#2）#3}}
\newcommand{\answercontinue}[2]{%
  \needspace{5\baselineskip}\subhdr{#1\quad 续解}%
  {\small\color{inkgray}本页接续：#2}\par\medskip}
% 延迟提供默认值，让旧讲义可在导言区保留自己的 \newcommand{\step}。
\AtBeginDocument{\providecommand{\step}[1]{\needspace{3\baselineskip}\par\medskip\textbf{#1}\par}}

% 前提旁注：只放在「需要理解的那一步」后面，给陈述 + 本题代入，不重证。
\newcommand{\aside}[1]{%
  \par\vspace{1pt}\noindent
  \hangindent=1.6em\hangafter=1
  {\small\color{inkgray}%
    \hspace*{0.6em}{\color{petal}$\triangleright$}\hspace{2pt}#1}%
  \par\vspace{2pt}}

%---------------------------------------------------------------- 表格
% 本机 MiKTeX 的 colortbl 与 array 版本不匹配（\rowcolor 会报错），
% 故表格统一用「上粗紫线 + 表头 + 中细线」的三线式，不依赖 colortbl。
\newcommand{\tabtop}{\par\vspace{3pt}{\color{violet}\hrule height 1.2pt}\vspace{1pt}}
\newcommand{\tabmid}{\par\vspace{1pt}{\color{petal}\hrule height 0.5pt}\vspace{1pt}}
\newcommand{\tabbot}{\par\vspace{3pt}{\color{rulepurple}\hrule height 0.5pt}\vspace{8pt}}
\newcommand{\thead}[1]{{\bfseries\color{violetdeep}#1}}

%---------------------------------------------------------------- 页眉页脚
\pagestyle{fancy}
\fancyhf{}
\renewcommand{\headrulewidth}{0.4pt}
\renewcommand{\headrule}{\hbox to\headwidth{\color{rulepurple}\leaders\hrule height \headrulewidth\hfill}}
\fancyhead[L]{\small\color{inkgray}\songti \doctag}
\fancyhead[R]{\small\color{inkgray}\songti \coursetag}
\fancyfoot[C]{\small\color{inkgray}\orn\enspace{\color{violet}\thepage}\enspace\orn}

% 每篇讲义在导言区设这两个值
\newcommand{\doctag}{}
\newcommand{\coursetag}{}
\newcommand{\settags}[2]{\renewcommand{\doctag}{#1}\renewcommand{\coursetag}{#2}}

%---------------------------------------------------------------- 数学记号
\newcommand{\vect}[1]{\boldsymbol{#1}}
\newcommand{\dd}{\,\mathrm{d}}
\newcommand{\ddx}{\frac{\mathrm{d}}{\mathrm{d}x}}
\addtolength{\jot}{2pt}



\setlength{\headheight}{20pt}
\renewcommand{\tabmid}{\noalign{\vskip 2pt{\color{petal}\hrule height 0.5pt}\vskip 3pt}}
\settags{EO \& MO · 作业答案卡}{光学（02）· 外教 Lecture 4}
\newcommand{\vdt}{\mathcal{V}}
\geometry{left=1.9cm,right=1.9cm,top=1.8cm,bottom=2.0cm}
\setlength{\parskip}{2.2pt}
\begin{document}

\begin{center}
{\LARGE\bfseries\color{violetdeep} Homework: Electro-optic \& Magneto-optic Effects}\\[5pt]
{\color{petal}\small \ding{95}\enspace Answer sheet (courseware p82)\enspace \ding{96}}\\[4pt]
{\small\color{inkgray}Handwrite in English;\ keep the four parts \textbf{Given} / \textbf{Formulas used} /
\textbf{calculation} / \textbf{Answers with units};\ sign your name in English.}
\end{center}
\vspace{-2pt}
{\color{petal}\hrule height 1pt}

\subhdr{Problem 1 \ ---\ Pockels cell phase retardation}

\textbf{Given}\quad KDP: $r=10.6\times10^{-12}\ \mathrm{m/V}$, $n=1.5$; crystal length
$L=3.0\times10^{-2}\ \mathrm{m}$, electrode spacing (aperture) $d=5.0\times10^{-3}\ \mathrm{m}$,
applied voltage $V=5.0\times10^{3}\ \mathrm{V}$; $\lambda=633\ \mathrm{nm}$. The field is across $d$ while the
light travels along $L$, so this is the \emph{transverse} configuration.

\textbf{Formulas used}\quad $E=V/d$;\quad
$\Delta\phi=\dfrac{2\pi}{\lambda}n^{3}rE\,L$ (the two eigen-polarizations differ by $n^3rE$);\quad
$V_\pi=(\lambda/2n^{3}r)\cdot(d/L)$;\quad $\Delta\phi=\pi V/V_\pi$.

\textbf{Calculation}\quad
\begin{equation*}
E=\frac{5.0\times10^{3}}{5.0\times10^{-3}}=1.0\times10^{6}\ \mathrm{V/m},
\end{equation*}
\begin{equation*}
\Delta\phi=\frac{2\pi}{633\times10^{-9}}\times1.5^{3}\times10.6\times10^{-12}\times1.0\times10^{6}\times3.0\times10^{-2}
\approx1.07\times10^{1}\ \mathrm{rad},
\end{equation*}
\begin{equation*}
V_\pi=\frac{633\times10^{-9}}{2\times1.5^{3}\times10.6\times10^{-12}}\times\frac{5.0\times10^{-3}}{3.0\times10^{-2}}
\approx1.47\times10^{3}\ \mathrm{V}.
\end{equation*}
Check:\ $V_\pi=\pi V/\Delta\phi=\pi\times5000/10.65\approx1.47\times10^{3}\ \mathrm{V}$.

\textbf{Answers}\quad (a) $\Delta\phi\approx1.1\times10^{1}\ \mathrm{rad}\approx610^{\circ}\ (\approx3.4\pi)$;\qquad
(b) $V_\pi\approx1.5\times10^{3}\ \mathrm{V}\ (\approx1.5\ \mathrm{kV})$.

\newpage
\subhdr{Problem 2 \ ---\ Faraday rotation and a double pass}

\textbf{Given}\quad $\lambda=589\ \mathrm{nm}$, $L=10\ \mathrm{cm}=0.10\ \mathrm{m}$, $B=0.5\ \mathrm{T}$,
$\vdt=20\ \mathrm{rad/(T\cdot m)}$; single pass in (a), double pass (back and forth) in (b).

\textbf{Formulas used}\quad $\theta=\vdt BL$;\quad
$\theta_{\text{round trip}}=2\vdt BL$ (the Faraday effect is \emph{non-reciprocal}).

\textbf{Calculation}\quad
\begin{equation*}
\text{(a)}\quad \theta=\vdt BL=20\times0.5\times0.10=1.0\ \mathrm{rad}\approx57.3^{\circ}.
\end{equation*}
\begin{equation*}
\text{(b)}\quad \theta_{\text{round trip}}=2\vdt BL=2.0\ \mathrm{rad}\approx115^{\circ}.
\end{equation*}

\textbf{Explanation}\quad The sense of rotation is fixed by the magnetic field, not by the direction of
propagation. On reflection $\mathbf k$ reverses but $\mathbf B$ does not, so the dielectric tensor and the two
circular refractive indices $n_\pm$ are unchanged; the relative phase $\frac{2\pi}{\lambda}(n_+-n_-)L$
accumulates a second time in the laboratory frame. The two passes therefore add, giving $2\theta$, whereas a
\emph{reciprocal} rotator (natural optical activity) would give $-\theta$ on the return pass and cancel to zero.
The doubling is not a coincidence: it is the signature of non-reciprocity.

\textbf{Answers}\quad (a) $\theta=1.0\ \mathrm{rad}=57.3^{\circ}$;\qquad
(b) $\theta_{\text{round trip}}=2.0\ \mathrm{rad}=115^{\circ}$, because the rotation sense is set by
$\mathbf B$, not by $\mathbf k$, so the return pass adds another $+\theta$.

\subhdr{Problem 3 \ ---\ Kerr effect versus Pockels effect}

\textbf{Given}\quad Crystal A: centrosymmetric; crystal B: non-centrosymmetric; $E=10^{6}\ \mathrm{V/m}$ for both.
Kerr data: $K=2\times10^{-14}\ \mathrm{m/V^{2}}$, $\lambda=633\ \mathrm{nm}$, $L=1.0\times10^{-2}\ \mathrm{m}$.

\textbf{Formulas used}\quad
\begin{itemize}
\item $\Delta n=K\lambda E^{2}$ (quadratic electro-optic effect);
\item $\Delta\phi=\dfrac{2\pi}{\lambda}\Delta n\,L$ (phase difference from birefringence);
\item symmetry rule: in a centrosymmetric medium $\Delta n(+E)=\Delta n(-E)$, so every odd-order
      term of $\Delta n(E)$ vanishes.
\end{itemize}

\textbf{Calculation}\quad
(a) Crystal A has inversion symmetry, so the linear term is forbidden --- it shows the \textbf{Kerr effect only}.
Crystal B allows the linear term --- it shows the \textbf{Pockels effect} (plus the Kerr effect, which is
hidden under the much larger linear term).

(b)
\begin{equation*}
\Delta n=K\lambda E^{2}=2\times10^{-14}\times633\times10^{-9}\times\left(10^{6}\right)^{2}
\approx1.27\times10^{-8}\quad(\text{dimensionless}).
\end{equation*}
(c)
\begin{equation*}
\Delta\phi=\frac{2\pi}{\lambda}\Delta n\,L
=\frac{2\pi}{633\times10^{-9}}\times1.27\times10^{-8}\times1.0\times10^{-2}
\approx1.26\times10^{-3}\ \mathrm{rad}.
\end{equation*}
Equivalent check (substituting $\Delta n$ first cancels $\lambda$):
$\Delta\phi=2\pi KLE^{2}\approx1.26\times10^{-3}\ \mathrm{rad}$.

\textbf{Answers}\quad (a) A: Kerr only; B: Pockels (linear) and Kerr;\qquad
(b) $\Delta n\approx1.3\times10^{-8}$;\qquad
(c) $\Delta\phi\approx1.3\times10^{-3}\ \mathrm{rad}\ (\approx0.07^{\circ})$.

\vspace{4pt}
{\color{petal}\hrule height 0.6pt}
{\small\color{inkgray}\textbf{Before handing in:} handwritten; four parts per problem; answers with units;
name signed in English; placed on the instructor's desk before the lecture starts.}

\end{document}
